\chapter*{摘\quad 要}
\addcontentsline{toc}{chapter}{摘要}

在传统操作系统课程中，学生需要先理解进程、虚拟内存、文件系统和设备驱动等机制，再修改内核代码，并通过编译、QEMU 测试或开发板实验验证结果。生成式 AI 能快速给出代码，却不会自动满足课程的教学要求。它可能在学生尚未形成设计思路时直接给出答案，也可能生成一份能够编译但学生无法解释的补丁。学生如果直接接受模型结果，便难以说明代码为什么正确、修改影响哪些模块、测试验证了什么，以及换一个环境后能否复现。教师也无法只看最终代码判断学生是否真正理解实验。

VeriSpecOSLab 将 AI 放进一套由学生主导的实验流程。学生可以先向问答助手讨论体系结构、内核组织、保护机制和目标硬件，再亲手填写系统设计、模块职责、跨边界接口和测试要求。平台以确定性程序检查规格结构、模块引用、文件范围和验证映射。评审助手读取学生已经写好的规格，指出设计遗漏和相互矛盾之处，同时保持项目文件不变。学生根据意见修改规格并用普通 Git 命令提交，设计决定由此始终来自学生。

AI 实现代码前，学生需要先提交当前规格。平台随后从这一提交创建独立的临时工作目录，并把目标模块允许修改的文件范围交给实现助手。助手交付内核实现，同时提出公开测试、契约测试、固定种子 fuzz、有限 trace 测试和本地 hidden tests。平台校验测试描述，读取 Git 实际记录的文件变化，再执行构建与全部相关测试。通过检查的实现形成一份带 Run-ID 和 Spec-Hash 的提交。调试、验证、问答和评审助手保持只读，学生可以清楚区分设计责任、实现帮助和最终验收。

平台统一管理编译、QEMU 启动、开发板运行和验证步骤。每次执行都会保存所用代码版本、实际命令、完整日志、输出文件和判定结果，因此学生能够回看一次实验究竟做了什么，教师也能够根据同一份记录复查。开发板记录由教学团队结合串口输出和板卡状态完成复核。课程资料可以导入本地知识库，问答助手在回答中给出来源、revision、content hash 和片段，便于学生回到教材或代码继续理解。

仓库提供一套包含完整源码的 RISC-V xv6 课程参考实现。Lab 1--10 依次覆盖最小工程、QEMU 启动、内存、进程、文件系统、完整 usertests、VisionFive 2 移植和报告提交。本文依据当前仓库的代码、文档、测试和历史记录，介绍 VeriSpecOSLab 的设计目标、系统组成、实现方法、验证结果、第三方来源、AI 协作记录和复现步骤。

\clearpage
